\chapter{Rust Patterns for Finance}
\label{app:rust}

This appendix documents the Rust patterns that recur across the
\crate{quant-lab} crates. Each pattern is illustrated with a minimal
example drawn from the actual codebase. The goal is to give the reader
a vocabulary for the design decisions made in the companion code and a
mental model for extending it.

\section{The Trait as Extension Point}

The single most important pattern in this book is the \emph{trait as
extension point}. The library defines a trait; the user implements it
for their specific model; the generic engine drives the computation.
This is how every pluggable component in \crate{quant-lab} works.

\begin{lstlisting}[language=Rust]
// quant-core/src/traits/labeler.rs
pub trait Labeler {
    type Input;
    type Label;

    fn label(&self, input: &Self::Input) -> Vec<Self::Label>;
}
\end{lstlisting}

A concrete implementation---the triple-barrier labeler from
\crate{quant-backtest}---supplies the method body:

\begin{lstlisting}[language=Rust]
pub struct TripleBarrier {
    pub take_profit: f64,
    pub stop_loss: f64,
    pub time_barrier: usize,
}

impl Labeler for TripleBarrier {
    type Input = (Vec<f64>, usize); // (prices, entry_index)
    type Label = i8;                // +1, -1, or 0

    fn label(&self, input: &Self::Input) -> Vec<Self::Label> {
        // ... barrier-touch logic ...
    }
}
\end{lstlisting}

\begin{keyinsight}
The trait defines the contract; the struct implements it; the generic
engine is oblivious to the concrete type. This is how Rust achieves
polymorphism without dynamic dispatch when performance matters
(monomorphisation) and with dynamic dispatch when flexibility matters
(\texttt{dyn Trait}).
\end{keyinsight}

\section{Pluggable Component Trait Catalogue}

The \crate{quant-core} traits module exposes the cross-cutting
interfaces used by the higher-level crates. Each trait corresponds to a
pluggable decision in a quantitative pipeline.

\begin{center}
\begin{tabular}{lp{8cm}}
\toprule
\textbf{Trait} & \textbf{Role} \\
\midrule
\func{CrossValidator} & Produces train/test index splits (k-fold, walk-forward, purged) \\
\func{Labeler} & Maps a price series to a label series (triple-barrier, fixed-horizon) \\
\func{BetSizer} & Converts a signal into a position size (Kelly, equal, fixed) \\
\func{SampleWeighter} & Assigns per-sample weights for training (return-based, time-decay) \\
\func{StructuralBreakDetector} & Locates regime-change indices in a series (CUSUM, Z-score) \\
\func{StochasticProcess} & Simulates a path $(S_0, \dots, S_T)$ (GBM, Brownian) \\
\func{OptionPricer} & Prices a vanilla option (Black-Scholes, Monte Carlo) \\
\func{Greeks} & Computes risk sensitivities ($\Delta, \Gamma, \nu, \Theta, \rho$) \\
\func{FactorModel} & Maps returns to factor exposures (PCA, Fama-French) \\
\func{ImpactModel} & Estimates market impact of an order (square-root, linear) \\
\func{OrderBookOps} & Order book operations (add, cancel, match, depth) \\
\bottomrule
\end{tabular}
\end{center}

Every trait in this table is used by at least one concrete struct in
the companion crate, and every concrete struct is used in at least one
example in \crate{quant-lib}.

\section{Generic Engines with Static Dispatch}

A backtest engine is generic over its labeler, cross-validator, and
bet sizer. The user composes these at construction time and the compiler
generates a specialised binary---no vtable overhead on the hot path.

\begin{lstlisting}[language=Rust]
// quant-backtest/src/generic.rs
pub struct GenericBacktest<L, CV, BS>
where
    L: Labeler,
    CV: CrossValidator,
    BS: BetSizer,
{
    labeler: L,
    validator: CV,
    sizer: BS,
    // ...
}

impl<L, CV, BS> GenericBacktest<L, CV, BS>
where
    L: Labeler,
    CV: CrossValidator,
    BS: BetSizer,
{
    pub fn run(&self, prices: &[f64]) -> BacktestResult {
        let labels = self.labeler.label(&(prices.to_vec(), 0));
        let splits = self.validator.splits(0, prices.len());
        // ... drive the loop, calling self.sizer.size(...) per bar ...
    }
}
\end{lstlisting}

The bound \texttt{where L: Labeler, CV: CrossValidator, BS: BetSizer}
is the entire contract. Adding a new labeler means writing an
\texttt{impl Labeler for MyLabeler}---no engine changes.

\section{The Builder Pattern for Composable Backtests}

When a struct has many optional parameters, a builder gives a fluent
API while keeping the internal representation compact.

\begin{lstlisting}[language=Rust]
// quant-backtest/src/generic.rs
pub struct BacktestBuilder {
    labeler: Option<Box<dyn Labeler<Input = (Vec<f64>, usize), Label = i8>>>,
    sizer: Option<Box<dyn BetSizer>>,
    // ... other pluggable components ...
}

impl BacktestBuilder {
    pub fn new() -> Self { /* all None */ }

    pub fn with_labeler<L: Labeler + 'static>(mut self, l: L) -> Self {
        self.labeler = Some(Box::new(l));
        self
    }

    pub fn with_sizer<BS: BetSizer + 'static>(mut self, bs: BS) -> Self {
        self.sizer = Some(Box::new(bs));
        self
    }

    pub fn build(self) -> Result<GenericBacktest<...>, BacktestError> {
        // validate required fields, return Err if any missing
    }
}
\end{lstlisting}

Usage reads as a pipeline:

\begin{lstlisting}[language=Rust]
let bt = BacktestBuilder::new()
    .with_labeler(TripleBarrier::new(0.02, 0.01, 5))
    .with_sizer(KellySizer::new(0.5))
    .build()?;
\end{lstlisting}

\marginnote{%
\textbf{Static vs.\ dynamic}\\[0.3em]
The engine uses static dispatch (generics) for the hot path, but the
builder uses \texttt{Box<dyn Trait>} for storage convenience. The
trade-off is a single vtable call per bar, negligible against the
per-bar computation.
}

\section{Newtype Wrappers for Domain Safety}

A bare \texttt{f64} can be a price, a return, a volatility, or a
probability---the compiler cannot tell them apart. A newtype wrapper
adds zero-cost type safety.

\begin{lstlisting}[language=Rust]
pub struct Price(f64);
pub struct LogReturn(f64);
pub struct Volatility(f64);

impl Price {
    pub fn log_return(&self, prev: &Price) -> LogReturn {
        LogReturn((self.0 / prev.0).ln())
    }
}
\end{lstlisting}

The compiler now refuses to add a \texttt{Price} to a
\texttt{Volatility}, even though both wrap an \texttt{f64}. The
\crate{quant-lab} crates use newtypes sparingly---mostly in public
APIs---to avoid ceremony in internal computation.

\section{Error Handling with \texttt{Result}}

Library code never panics. Every fallible operation returns
\texttt{Result<T, E>}, with the error type implementing
\texttt{std::error::Error}. The crate-level error is an enum:

\begin{lstlisting}[language=Rust]
// quant-core/src/error.rs
#[derive(Debug, thiserror::Error)]
pub enum QuantError {
    #[error("insufficient data: required {required}, got {actual}")]
    InsufficientData { required: usize, actual: usize },

    #[error("singular matrix: {0}")]
    SingularMatrix(String),

    #[error("optimisation failed to converge")]
    NonConvergence,

    #[error(transparent)]
    Io(#[from] std::io::Error),
}
\end{lstlisting}

The \texttt{InsufficientData} variant appears in the ADF test
(too few observations for the lag order), \texttt{SingularMatrix} in
OLS regression (collinear regressors), and \texttt{NonConvergence} in
implied-volatility inversion. Callers use \texttt{?} to propagate:

\begin{lstlisting}[language=Rust]
fn sharpe(returns: &[f64], rf: f64) -> Result<f64, QuantError> {
    let m = mean(returns);
    let s = std_dev(returns).ok_or(QuantError::InsufficientData {
        required: 2,
        actual: returns.len(),
    })?;
    Ok((m - rf) / s * 252.0_f64.sqrt())
}
\end{lstlisting}

\section{Iterators and Lazy Evaluation}

Most transformations are expressed as iterator adaptors, avoiding
intermediate allocations:

\begin{lstlisting}[language=Rust]
let mean: f64 = returns.iter().sum::<f64>() / returns.len() as f64;
let var: f64 = returns.iter()
    .map(|r| (r - mean).powi(2))
    .sum::<f64>() / (returns.len() - 1) as f64;
\end{lstlisting}

For rolling windows, the \func{RollingWindow} trait returns an
iterator of slices, letting the caller fold, map, or collect without
the library dictating the output shape.

\section{Feature Flags for Module Gating}

The \crate{quant-lib} facade crate exposes every submodule behind a
feature flag, so users pay only for what they use:

\begin{lstlisting}
[features]
default = ["all"]
all = ["core", "timeseries", "vol", "stochastic",
        "options", "portfolio", "factors",
        "microstructure", "backtest"]
\end{lstlisting}

In Rust source, each module is conditionally compiled:

\begin{lstlisting}[language=Rust]
#[cfg(feature = "portfolio")]
pub mod portfolio;
\end{lstlisting}

A user who needs only option pricing can depend on
\crate{quant-lib} with \texttt{default-features = false, features =
["options"]}, and the rest of the crate is not compiled.

\section{The Facade Crate Pattern}

\crate{quant-lib} is a \emph{facade}: it contains \texttt{pub mod}
and \texttt{pub use} declarations, zero new logic. It re-exports the
public APIs of \crate{quant-core}, \crate{quant-timeseries},
\crate{quant-vol}, \dots\ under a single, versioned namespace.

\begin{lstlisting}[language=Rust]
// quant-lib/src/lib.rs
pub use quant_core;
pub use quant_timeseries;
pub use quant_vol;
// ...

pub mod prelude {
    pub use crate::core::*;
    pub use crate::timeseries::*;
    // ...
}
\end{lstlisting}

Users write \texttt{use quant\_lib::prelude::*;} and get the entire
public surface. Upgrading a single underlying crate (e.g.,
\crate{quant-options} for a new pricing model) is a one-line change
in \crate{quant-lib/Cargo.toml}; downstream code is unaffected.

\section{Deterministic Randomness}

Every stochastic example in this book uses a seeded RNG to be
reproducible:

\begin{lstlisting}[language=Rust]
use quant_lib::prelude::*;
let mut rng = XorShift64::new(42);
let path = gbm(100.0, 0.05, 0.20, 1.0, 252, &mut rng)?;
\end{lstlisting}

The \func{Rng} trait abstracts the RNG so a Monte Carlo engine can
swap \texttt{XorShift64} for a cryptographic generator without code
changes. Tests use a fixed seed to make results bit-for-bit
reproducible across machines and CI runs.

\section{Testing Strategy}

Every crate ships three kinds of tests:

\begin{itemize}
    \item \textbf{Unit tests} (in \texttt{src/}): small, fast, test one
      function in isolation.
    \item \textbf{Integration tests} (in \texttt{tests/}): exercise the
      public API end-to-end.
    \item \textbf{Examples} (in \texttt{examples/}): runnable
      narratives; each chapter's companion code is an example.
\end{itemize}

Examples are the primary documentation: \texttt{cargo run -p quant-lib
-example 08\_options\_greeks} reproduces the numbers printed in the
chapter.

\begin{keyinsight}
The trait-based architecture makes the library testable at the
composition boundary: a test that passes with a
\texttt{FixedHorizonLabeler} will also pass with a
\texttt{TripleBarrierLabeler} because both satisfy the same
\texttt{Labeler} contract. Swapping implementations in tests is a
one-line change.
\end{keyinsight}